% GENERATED FILE. Edit canonical lessons, then run npm run generate:books.
% canonical-source-hash: fnv1a64:22a549265d659d3c
% canonical-lessons: ZH-C66-tianqi, ZH-W66-qi-air, ZH-C66-yu, ZH-W66-yu-rain, ZH-C66-xiayu, ZH-C66-xue, ZH-C66-feng

\chapter{The Weather}
\label{ch:zh-the-weather}
\hlchaptermodality{66}

\begin{chapteropening}
\textbf{By the end of this chapter:} \emph{I can say what the weather is like: rain, snow and wind.}

Complete the last lesson of chapter 66: I can say what the weather is like: rain, snow and wind.
\end{chapteropening}

\section[\texorpdfstring{\zh{天气} (tiānqì)}{天气 (tiānqì)}]{\zh{天气} (tiānqì) — weather}

\label{lesson:ZH-C66-tianqi}



\begin{quote}
Before the new one: say the Mandarin for to open, then the Mandarin for to shut.
\end{quote}

\subsection*{You'll want to know: \zh{天气}}
\textbf{\zh{天气}} — \emph{tiānqì} — \textquotedblleft{}weather\textquotedblright{}.

\begin{quote}
\textbf{\zh{今天天气很好}} \emph{jīntiān tiānqì hěn hǎo} — The weather is lovely today.
\end{quote}

The new character, in pieces: \zh{气} — a short fall, two short bars, and a long stroke that turns, slants out and hooks. The next lesson puts it in your hand.

\begin{cousinweb}
\textbf{\zh{天气}} (\emph{tiānqì}) is built from \zh{天} 'sky' and \zh{气} 'breath', both inherited from Middle Chinese, \emph{then} and \emph{khj+jH} (Guangyun, in Baxter's transcription): the breath of the sky.
\end{cousinweb}

\begin{grammarlens}[title={What you've built}]
One more word for saying how you are.
\end{grammarlens}

\subsection*{Your turn}
\begin{itemize}
\raggedright
  \item \emph{Say it:} \emph{tiānqì}
  \item \emph{Say it:} \emph{tiānqì}, once more
  \item \emph{Say it:} \emph{jīntiān tiānqì hěn hǎo} — The weather is lovely today.
  \item \emph{Recall:} say the Mandarin for to open, then the Mandarin for to shut, then say \emph{tiānqì} again
\end{itemize}

\begin{tcolorbox}[breakable,colback=teal!4,colframe=teal!35!black,title={Before you move on}]
What does \zh{天气} mean? (\textquotedblleft{}Weather\textquotedblright{}.) Say it once more.
\end{tcolorbox}
\section[\texorpdfstring{\zh{气} (qì)}{气 (qì)}]{\zh{气} — a fall, two bars, a hook}

\label{lesson:ZH-W66-qi-air}



\begin{quote}
The word first, then the character inside it.

\begin{itemize}
\raggedright
  \item \emph{Recall:} say \emph{tiānqì} — \textbf{R1}, one lesson back, before it decays
\end{itemize}

\begin{quote}
\zh{天} \zh{手}
\end{quote}

\emph{Read it:} each one aloud, with nothing said first — \emph{tiān}, \emph{shǒu}
\end{quote}

\subsection*{Script — a fall, two bars, a hook}
\hlblockfigure{\detokenize{figures/ZH-W66-qi-air-filmstrip.pdf}}{How \zh{气} is written, stroke by stroke}

\begin{quote}
\zh{气}
\end{quote}

\textbf{One piece, four strokes.}

\textbf{A short left-falling stroke} at the top left.

\textbf{Two short horizontals} under it, the second a little lower and shorter.

\textbf{A long horizontal that turns down}, slants out to the right and hooks up at the end.

\subsection*{Your turn}
\begin{itemize}
\raggedright
  \item \emph{Write it:} \zh{气}, all four, with the model in front of you
  \item \emph{Cover:} write it again from memory, and count the pieces, not the strokes
\end{itemize}

\begin{tcolorbox}[breakable,colback=teal!4,colframe=teal!35!black,title={Before you move on}]
How many strokes in \zh{气}? (\textbf{Four}.) How many short bars? (\textbf{Two}.) How does the last stroke end? (\textbf{It slants out and hooks up}.)

Source: \href{https://raw.githubusercontent.com/chanind/hanzi-writer-data/68d10a4b21150cae5e1ebbd223eed289cf32d90c/data/\%E6\%B0\%94.json}{Hanzi Writer Data \zh{气}, pinned stroke paths and medians}.
\end{tcolorbox}
\section[\texorpdfstring{\zh{雨} (yǔ)}{雨 (yǔ)}]{\zh{雨} (yǔ) — rain}

\label{lesson:ZH-C66-yu}



\begin{quote}
Before the new one: say the Mandarin for to shut, then the Mandarin for weather.
\end{quote}

\subsection*{You'll want to know: \zh{雨}}
\textbf{\zh{雨}} — \emph{yǔ} — \textquotedblleft{}rain\textquotedblright{}.

\begin{quote}
\textbf{\zh{请问明天有雨吗}} \emph{qǐngwèn míngtiān yǒu yǔ ma} — Excuse me, will it rain tomorrow?
\end{quote}

The new character, in pieces: \zh{雨} — a lid, a frame with a line down its middle, and four drops inside. The next lesson puts it in your hand.

\begin{cousinweb}
\textbf{\zh{雨}} (\emph{yǔ}) is inherited from Middle Chinese \emph{hjuX} (Guangyun, in Baxter's transcription), which the dictionary explains, through an older book, as water coming down from the clouds.
\end{cousinweb}

\begin{grammarlens}[title={What you've built}]
One more word for saying how you are.
\end{grammarlens}

\subsection*{Your turn}
\begin{itemize}
\raggedright
  \item \emph{Say it:} \emph{yǔ}
  \item \emph{Say it:} \emph{yǔ}, once more
  \item \emph{Say it:} \emph{qǐngwèn míngtiān yǒu yǔ ma} — Excuse me, will it rain tomorrow?
  \item \emph{Recall:} say the Mandarin for to shut, then the Mandarin for weather, then say \emph{yǔ} again
\end{itemize}

\begin{tcolorbox}[breakable,colback=teal!4,colframe=teal!35!black,title={Before you move on}]
What does \zh{雨} mean? (\textquotedblleft{}Rain\textquotedblright{}.) Say it once more.
\end{tcolorbox}
\section[\texorpdfstring{\zh{雨} (yǔ)}{雨 (yǔ)}]{\zh{雨} — a lid, a frame and four drops}

\label{lesson:ZH-W66-yu-rain}



\begin{quote}
The word first, then the character inside it.

\begin{itemize}
\raggedright
  \item \emph{Recall:} say \emph{yǔ} — \textbf{R1}, one lesson back, before it decays
\end{itemize}

\begin{quote}
\zh{天} \zh{气}
\end{quote}

\emph{Read it:} each one aloud, with nothing said first — \emph{tiān}, \emph{qì}
\end{quote}

\subsection*{Script — a lid, a frame and four drops}
\hlblockfigure{\detokenize{figures/ZH-W66-yu-rain-filmstrip.pdf}}{How \zh{雨} is written, stroke by stroke}

\begin{quote}
\zh{雨}
\end{quote}

\textbf{One piece, built from the top down.}

\textbf{A horizontal} across the top, the lid.

\textbf{A frame, two strokes}: a short vertical down the left, then a horizontal that turns down the right side and hooks to the left.

\textbf{A vertical} down the middle from the lid.

\textbf{Four dots}, two on the left of the middle line, then two on the right.

1 + 2 + 1 + 4 = eight.

\subsection*{Your turn}
\begin{itemize}
\raggedright
  \item \emph{Write it:} \zh{雨}, all eight, with the model in front of you
  \item \emph{Cover:} write it again from memory, and count the pieces, not the strokes
\end{itemize}

\begin{tcolorbox}[breakable,colback=teal!4,colframe=teal!35!black,title={Before you move on}]
How many strokes in \zh{雨}? (\textbf{Eight}.) What comes last? (\textbf{The four dots, the left pair first}.) Which stroke hooks? (\textbf{The right side of the frame}.)

Source: \href{https://raw.githubusercontent.com/chanind/hanzi-writer-data/68d10a4b21150cae5e1ebbd223eed289cf32d90c/data/\%E9\%9B\%A8.json}{Hanzi Writer Data \zh{雨}, pinned stroke paths and medians}.
\end{tcolorbox}
\section[xiàyǔ]{xiàyǔ — to rain}

\label{lesson:ZH-C66-xiayu}



\begin{quote}
Before the new one: say the Mandarin for weather, then the Mandarin for rain.
\end{quote}

\subsection*{You'll want to know: xiàyǔ}
\emph{xiàyǔ} — \textquotedblleft{}to rain\textquotedblright{}.

\begin{quote}
\emph{xiàyǔ wǒmen bù dǎqiú} — It's raining, so we're not playing ball.
\end{quote}

Its characters come later; for now it is a word for the ear.

\begin{cousinweb}
\emph{xiàyǔ} is built from xià 'to come down' and \zh{雨} 'rain', both inherited from Middle Chinese, \emph{haeX} and \emph{hjuX} (Guangyun, in Baxter's transcription): rain comes down.
\end{cousinweb}

\begin{grammarlens}[title={What you've built}]
One more word for saying how you are.
\end{grammarlens}

\subsection*{Your turn}
\begin{itemize}
\raggedright
  \item \emph{Say it:} \emph{xiàyǔ}
  \item \emph{Say it:} \emph{xiàyǔ}, once more
  \item \emph{Say it:} \emph{xiàyǔ wǒmen bù dǎqiú} — It's raining, so we're not playing ball.
  \item \emph{Recall:} say the Mandarin for weather, then the Mandarin for rain, then say \emph{xiàyǔ} again
\end{itemize}

\begin{tcolorbox}[breakable,colback=teal!4,colframe=teal!35!black,title={Before you move on}]
What does xiàyǔ mean? (\textquotedblleft{}To rain\textquotedblright{}.) Say it once more.
\end{tcolorbox}
\section[xuě]{xuě — snow}

\label{lesson:ZH-C66-xue}



\begin{quote}
Before the new one: say the Mandarin for rain, then the Mandarin for to rain.
\end{quote}

\subsection*{You'll want to know: xuě}
\emph{xuě} — \textquotedblleft{}snow\textquotedblright{}.

\begin{quote}
\emph{xuě hěn dà wǒ de ěrduo hěn lěng} — It's snowing hard and my ears are freezing.
\end{quote}

Its characters come later; for now it is a word for the ear.

\begin{cousinweb}
\emph{xuě} is inherited from Middle Chinese \emph{sjwet} (Guangyun, in Baxter's transcription), glossed 'frozen rain'.
\end{cousinweb}

\begin{grammarlens}[title={What you've built}]
One more word for saying how you are.
\end{grammarlens}

\subsection*{Your turn}
\begin{itemize}
\raggedright
  \item \emph{Say it:} \emph{xuě}
  \item \emph{Say it:} \emph{xuě}, once more
  \item \emph{Say it:} \emph{xuě hěn dà wǒ de ěrduo hěn lěng} — It's snowing hard and my ears are freezing.
  \item \emph{Recall:} say the Mandarin for rain, then the Mandarin for to rain, then say \emph{xuě} again
\end{itemize}

\begin{tcolorbox}[breakable,colback=teal!4,colframe=teal!35!black,title={Before you move on}]
What does xuě mean? (\textquotedblleft{}Snow\textquotedblright{}.) Say it once more.
\end{tcolorbox}
\section[fēng]{fēng — wind}

\label{lesson:ZH-C66-feng}



\begin{quote}
Before the new one: say the Mandarin for to rain, then the Mandarin for snow.
\end{quote}

\subsection*{You'll want to know: fēng}
\emph{fēng} — \textquotedblleft{}wind\textquotedblright{}.

\begin{quote}
\emph{tiānqì bù hǎo fēng hěn dà} — The weather's bad: the wind is strong.
\end{quote}

Its characters come later; for now it is a word for the ear.

\begin{cousinweb}
\emph{fēng} is inherited from Middle Chinese \emph{pjuwng} (Guangyun, in Baxter's transcription); the dictionary quotes an old book calling the wind the messenger of sky and earth.
\end{cousinweb}

\begin{grammarlens}[title={What you've built}]
One more word for saying how you are.
\end{grammarlens}

\subsection*{Your turn}
\begin{itemize}
\raggedright
  \item \emph{Say it:} \emph{fēng}
  \item \emph{Say it:} \emph{fēng}, once more
  \item \emph{Say it:} \emph{tiānqì bù hǎo fēng hěn dà} — The weather's bad: the wind is strong.
  \item \emph{Recall:} say the Mandarin for to rain, then the Mandarin for snow, then say \emph{fēng} again
\end{itemize}

\begin{tcolorbox}[breakable,colback=teal!4,colframe=teal!35!black,title={Before you move on}]
What does fēng mean? (\textquotedblleft{}Wind\textquotedblright{}.) Say it once more.
\end{tcolorbox}
